%%%%%% -*- Coding: utf-8-unix; Mode: latex

\documentclass[english]{projectreport}

\usepackage[utf8]{inputenc}
\usepackage{url}
\usepackage{subfigure}
\usepackage{tabularx}
\usepackage{ragged2e}
\usepackage{booktabs}
\usepackage{multirow}
\usepackage{grffile}
\usepackage{amsmath}

\usepackage[bookmarks=false]{hyperref}
\hypersetup{colorlinks,
  linkcolor=blue,
  citecolor=blue,
  urlcolor=blue}

\usepackage{indentfirst}
\usepackage{caption}
\usepackage{listings}
\usepackage[ruled,linesnumbered,lined]{algorithm2e}
\usepackage[svgnames]{xcolor}
\usepackage{inconsolata}
\usepackage{fancyhdr}

\usepackage[
style=numeric,
sorting=nyt,
isbn=false,
doi=true,
url=true,
backref=false,
backrefstyle=none,
maxnames=10,
giveninits=true,
abbreviate=true,
defernumbers=false,
backend=biber]{biblatex}
\addbibresource{bibliography.bib}

\lstset{
    language=Python,
    basicstyle=\ttfamily\footnotesize,
    backgroundcolor=\color{gray!5},
    commentstyle=\color{Green},
    keywordstyle=\color{Red},
    stringstyle=\color{Blue},
    numberstyle=\tiny\color{Black},
    escapeinside=`',
    frame=single,
    tabsize=2,
    rulecolor=\color{black!30},
    title=\lstname,
    breaklines=true,
    breakatwhitespace=true,
    framextopmargin=2pt,
    framexbottommargin=2pt,
    extendedchars=false,
    captionpos=b,
    abovecaptionskip=5pt,
    keepspaces=true,
    numbers=left,
    numbersep=5pt,
    showspaces=false,
    showstringspaces=false,
    showtabs=false,
    tabsize=2
}

\SetAlgorithmName{\LangAlgorithm}{\LangAlgorithmRef}{\LangListOfAlgorithms}
\renewcommand{\lstlistingname}{\LangListing}
\renewcommand\lstlistlistingname{\LangListOfListings}

\newcolumntype{Y}{>{\small\centering\arraybackslash}X}

\captionsetup[figure]{skip=5pt,position=bottom}
\captionsetup[table]{skip=5pt,position=top}


%%%%%%%%%%%%%%%%%%%%%%%%%%%%%%%%%%%%%%%%%%%%%%%%%%%%%%%%%%%%%%%%%%%%%%%%%%%%%%%
\author{Dawid Mularczyk$^1$}

\title{Multi-scale Agent-Based Simulation of Viral Pandemic Spread}

\date{
    {\small $^1$Computer Science, AGH University of Science and Technology \\
    \texttt{muldaw@student.agh.edu.pl}\\[2ex]%
    \the\year}
}

\pagestyle{fancy}
\fancyhf{}
\fancyhead[LE,RO]{\thepage}

%%%%%%%%%%%%%%%%%%%%%%%%%%%%%%%%%%%%%%%%%%%%%%%%%%%%%%%%%%%%%%%%%%%%%%%%%%%%%%%
\begin{document}

\maketitle
\thispagestyle{empty}

\begin{abstract}
\noindent
This document is a comprehensive report for a project focused on the design
and implementation of a multi-scale Agent-Based Model (ABM) simulating the
spread of a viral pandemic in an urban population. The project is carried out
in three stages: (1)~model design and initial implementation, (2)~full
implementation with preliminary verification and potential revisions,
(3)~proper experimental research and final project. The simulation platform
is built in Python using the Mesa framework. It implements a SEIRD model with
a continuous viral load variable and three transmission mechanisms: aerosol,
droplet, and indirect contact through contaminated surfaces (fomites).

\noindent\textbf{Keywords:} agent-based modeling, computational epidemiology,
SEIRD model, disease transmission, Mesa, Python, fomites
\end{abstract}

%%%%%%%%%%%%%%%%%%%%%%%%%%%%%%%%%%%%%%%%%%%%%%%%%%%%%%%%%%%%%%%%%%%%%%%%%%%%%%%
\section{\SectionTitleIntroduction}
\label{sec:introduction}

Understanding and predicting the dynamics of infectious disease spread is one
of the central challenges in modern computational epidemiology and public
health policy planning~\cite{anderson1992infectious}. Global health crises ---
from the 1918 Spanish flu, through the West African Ebola outbreak, to the
COVID-19 pandemic --- have demonstrated that computer simulation models are
indispensable decision-support tools: they allow testing intervention
scenarios, forecasting healthcare system load, and evaluating the effectiveness
of preventive measures before their real-world deployment~\cite{tracy2018agent}.

\subsection{Classical Compartmental Models}
\label{sec:compartmental}

Historically, the dominant approach has been compartmental models based on
systems of nonlinear ordinary differential equations (ODEs). The classical
SIR model~\cite{kermack1927sir} partitions a population into three groups:
Susceptible, Infectious, and Recovered. The transition dynamics are given by:

\begin{equation}
  \frac{dS}{dt} = -\beta \frac{SI}{N}, \quad
  \frac{dI}{dt} = \beta \frac{SI}{N} - \gamma I, \quad
  \frac{dR}{dt} = \gamma I
  \label{eq:sir}
\end{equation}

\noindent where $\beta$ is the transmission rate, $\gamma$ the recovery rate,
and $N$ the total population size~\cite{brauer2008compartmental}.

Despite high computational efficiency, the fundamental limitation of ODE
models is the homogeneous mixing assumption: every individual is equally
likely to contact any other. This ignores spatial structure, social clusters,
and individual behaviour~\cite{eubank2004modelling}.

\subsection{Agent-Based Modeling as a Response to ODE Limitations}
\label{sec:abm-motivation}

Agent-based modeling (ABM) represents a population as a collection of
autonomous individuals, each with their own set of attributes, internal state,
and behavioural rules~\cite{bonabeau2002agent}. Agents interact directly with
one another and with the environment in a defined topological space.
Microscopic interactions give rise to emergent macroscopic phenomena such as
epidemic waves, localised outbreak clusters, and superspreading
events~\cite{willem2017lessons}.

Key advantages of the agent-based approach for epidemic modeling:

\begin{itemize}
  \item explicit representation of population heterogeneity (age, behaviour,
        immunity),
  \item natural modeling of social structures (household, workplace, school),
  \item simultaneous simulation of multiple transmission routes (airborne,
        touch, fomites),
  \item emergent phenomena that deterministic ODE models cannot capture.
\end{itemize}

\subsection{Project Aim and Scope}
\label{sec:goal}

The aim of this project is to incrementally design and implement a multi-scale,
hybrid agent-based simulation of a viral pandemic in an urban environment. The
platform integrates compartmental mechanics (SEIRD model) with a full agent
architecture, enabling the analysis of individual behaviour, urban spatial
structure, and diverse epidemiological interventions on infection dynamics.
The project is structured into three stages:

\begin{enumerate}
  \item \textbf{Stage~1 --- Model and initial implementation:} architecture
        design, agent and environment specification, working SEIRD prototype.
  \item \textbf{Stage~2 --- Implementation and preliminary verification:} full
        implementation of transmission mechanics, parameter calibration, model
        correctness testing.
  \item \textbf{Stage~3 --- Proper experiment and final project:} experimental
        scenarios, results analysis, final conclusions.
\end{enumerate}

%%%%%%%%%%%%%%%%%%%%%%%%%%%%%%%%%%%%%%%%%%%%%%%%%%%%%%%%%%%%%%%%%%%%%%%%%%%%%%%
\section{Stage 1: Model and Initial Implementation}
\label{sec:stage1}

Stage~1 covers the complete design of the epidemiological model, specification
of the agent and environment architecture, selection of the technology stack,
and implementation of a working prototype.

\subsection{Epidemiological Model: SEIRD}
\label{sec:seird}

Each agent in the system transitions through an extended set of clinical
states --- the SEIRD model --- described by:

\begin{equation}
  S \;\xrightarrow{\lambda(t)}\; E \;\xrightarrow{\sigma}\; I
  \;\xrightarrow{\gamma}\; R, \qquad I \;\xrightarrow{\mu}\; D
  \label{eq:seird}
\end{equation}

\noindent where $\lambda(t)$ is the time-varying force of infection (dependent
on the agent's contacts), $\sigma$ the rate of progression from the incubation
period to the infectious phase, $\gamma$ the recovery rate, and $\mu$ the
disease-induced mortality rate.

The model states are:
\begin{itemize}
  \item \textbf{S} (\emph{Susceptible}) --- healthy agent with no immunity,
  \item \textbf{E} (\emph{Exposed}) --- agent in the incubation period;
        infected but not yet infectious (mobile, unaware of infection),
  \item \textbf{I} (\emph{Infectious}) --- actively infectious agent shedding
        pathogen,
  \item \textbf{R} (\emph{Recovered}) --- recovered agent with acquired
        immunity,
  \item \textbf{D} (\emph{Dead}) --- agent that did not survive the infection.
\end{itemize}

A key realism element is a continuous \emph{viral load} --- an internal agent
variable that rises toward the end of phase~E, peaks at the onset of phase~I,
and declines as the immune response progresses. Viral load directly scales the
probability of infecting neighbouring agents at each interaction step.

Stage durations are not fixed constants; they are sampled from a Gamma
distribution to capture natural biological variance~\cite{lloyd2001realistic}.

\subsection{Base Pathogen Parameters}
\label{sec:pathogen-params}

The modelled pathogen is a hypothetical aggressive respiratory virus with
characteristics intermediate between pandemic influenza and SARS-CoV-2.
Table~\ref{tab:pathogen} lists the initial parameter values, which will be
calibrated in Stage~2.

\begin{table}[!htbp]
\centering
\caption{Epidemiological parameters of the modelled pathogen (initial values)}
\begin{tabularx}{\columnwidth}{@{}lYY@{}} \toprule
  \textbf{Parameter} & \textbf{Value} & \textbf{Source / basis} \\ \midrule
  Basic reproduction number $R_0$ & 2.5\,--\,3.5 & Pandemic flu / COVID-19 \\
  Incubation period (phase E) & Gamma($\mu$=5d, $\sigma$=2d) & \cite{zhang2020covid} \\
  Infectious duration (phase I) & Gamma($\mu$=8d, $\sigma$=3d) & \cite{lloyd2001realistic} \\
  Case fatality rate (CFR) & 1.5\%\,--\,2.5\% (age-dependent) & \\
  Transmission routes & Aerosol, droplets, fomites & \\
  Post-infection immunity & 6\,--\,12 months & \\
  \bottomrule
\end{tabularx}
\label{tab:pathogen}
\end{table}

\subsection{Human Agent Attributes}
\label{sec:agent-attributes}

Each agent in the system (\texttt{HumanAgent}) is an instance of a class
aggregating three groups of attributes.

\subsubsection{Demographic and Sociological Attributes}
\label{sec:demo-attrs}

\begin{itemize}
  \item \textbf{Age} --- sampled from an empirical demographic pyramid;
        determines assignment to a functional group (child, adult, senior)
        and mobility patterns,
  \item \textbf{Household ID} --- agents sharing a household occupy the same
        space at night; household clusters are a primary transmission
        environment,
  \item \textbf{Workplace / school} --- an assigned POI node to which the
        agent travels regularly according to their daily schedule.
\end{itemize}

\subsubsection{Clinical and Immunological Attributes}
\label{sec:clinical-attrs}

\begin{itemize}
  \item \textbf{Disease state} --- current SEIRD compartment,
  \item \textbf{Viral load} --- continuous variable representing pathogen
        titre (used to compute transmission risk at each interaction),
  \item \textbf{Immunity profile} --- baseline biological susceptibility,
        modified by age and vaccination history; upon entering state~R the
        agent gains a temporary or permanent antibody pool,
  \item \textbf{Vaccination status} --- reduces infection probability and
        drastically lowers the risk of disease-induced death.
\end{itemize}

\subsubsection{Behavioural Attributes}
\label{sec:behavioral-attrs}

\begin{itemize}
  \item \textbf{Daily schedule} --- defines at which simulation hours the
        agent leaves home, travels to work, visits shops, etc.,
  \item \textbf{Hygiene score} ($\Theta_H$) --- personalised hand-washing
        frequency coefficient; actively reduces pathogen titre on hands,
        lowering the risk of self-inoculation and surface contamination,
  \item \textbf{Contact rate} --- heterogeneity in the number of daily
        interactions; agents with a high contact rate combined with a high
        viral load naturally emerge as superspreaders.
\end{itemize}

Table~\ref{tab:agent-attrs} summarises all planned attributes of the
\texttt{HumanAgent} class.

\begin{table}[!htbp]
\centering
\caption{Key attributes of the \texttt{HumanAgent} class}
\begin{tabularx}{\columnwidth}{@{}lYY@{}} \toprule
  \textbf{Attribute} & \textbf{Type / range} & \textbf{Role in model} \\ \midrule
  \texttt{age} & int, [0--100] & Mobility, immunity, mortality \\
  \texttt{household\_id} & int & Household clusters, nocturnal transmission \\
  \texttt{workplace\_id} & int / None & Schedule, POI node \\
  \texttt{state} & enum SEIRD & Agent clinical state \\
  \texttt{viral\_load} & float [0--1] & Pathogen emission strength \\
  \texttt{immunity} & float [0--1] & Susceptibility to infection \\
  \texttt{vaccinated} & bool & Transmission and mortality reduction \\
  \texttt{hygiene\_score} & float [0--1] & Hand washing, fomite reduction \\
  \texttt{contact\_rate} & float & Interactions per day \\
  \texttt{daily\_schedule} & list of POI & Sequence of nodes during the day \\
  \bottomrule
\end{tabularx}
\label{tab:agent-attrs}
\end{table}

\subsection{Environment Architecture: Space and POI Nodes}
\label{sec:environment}

The simulation environment is organised as a heterogeneous graph
$G = (V, E)$, where $V$ is the set of physical locations (Points of Interest,
POI) and $E$ the simplified transit paths connecting them~\cite{eubank2004modelling}.
Graph nodes are divided into semantic categories:

\begin{itemize}
  \item \textbf{Households} --- nocturnal and morning space,
  \item \textbf{Workplaces and offices} --- daily destination nodes for adults,
  \item \textbf{Educational facilities} --- schools and universities for
        children and young adults,
  \item \textbf{Commercial spaces} --- supermarkets and shops,
  \item \textbf{Recreational venues} --- parks, restaurants,
  \item \textbf{Healthcare facilities} --- clinics, hospitals.
\end{itemize}

Each POI node carries environmental attributes that affect aerosol dynamics:
capacity, floor area, ventilation coefficient, and temperature. Poorly
ventilated, densely occupied nodes (e.g.\ a classroom, a metro carriage)
dramatically increase the aerosol transmission risk compared to open spaces.

\subsection{Fomites: Passive Environmental Agents}
\label{sec:fomites}

A key element of the project is the modeling of indirect transmission via
contaminated surfaces --- fomites~\cite{jones2009fomites}. Selected high-risk
POI nodes (door handles, shared desks, handrails) contain objects of class
\texttt{FomiteAgent} --- passive agents with their own state:

\begin{itemize}
  \item \textbf{Contamination level} --- current pathogen titre on the surface,
  \item \textbf{Decay coefficient} ($\mu_F$) --- exponential decline of titre
        with each simulation clock tick.
\end{itemize}

The virus exchange mechanics between agents and fomites are governed by a set
of indirect transmission parameters (Table~\ref{tab:fomite-params}):

\begin{table}[!htbp]
\centering
\caption{Fomite transmission parameters}
\begin{tabularx}{\columnwidth}{@{}lYY@{}} \toprule
  \textbf{Symbol} & \textbf{Definition} & \textbf{Unit} \\ \midrule
  $\alpha$ & Pathogen shedding rate of an infectious host & virus copies / h \\
  $\phi$ & Fraction of shed load deposited on hands & [0--1] \\
  $\rho_T$ & Hand-to-surface contact frequency & contacts / h \\
  $\tau_{H \to F}$ & Transfer efficiency from hand to fomite & [0--1] \\
  $\tau_{F \to H}$ & Transfer efficiency from fomite to hand & [0--1] \\
  $\mu_H$ & Viral inactivation rate on hands & log decay / h \\
  $\mu_F$ & Viral inactivation rate on surface & log decay / h \\
  $\Theta_H$ & Hand-washing score (hygiene) & load reduction \\
  \bottomrule
\end{tabularx}
\label{tab:fomite-params}
\end{table}

\subsection{Direct Transmission Mechanics: Aerosols and Droplets}
\label{sec:aerosol}

Direct transmission is modelled as a spatially propagating physical event.
When an infectious agent sneezes or coughs, a spherical zone of elevated
aerosol concentration is generated around them. The probability that a
susceptible agent within this zone becomes infected is~\cite{nicas2004toward}:

\begin{equation}
  P_{\mathrm{inf}}(d, t) = 1 - \exp\!\left(
    -\frac{V_L \cdot f(d) \cdot \Delta t}{D_{\min}}
  \right)
  \label{eq:aerosol}
\end{equation}

\noindent where $V_L$ is the emitter's viral load, $f(d)$ a geometric decay
function with distance $d$ (gravitational settling of large droplets),
$\Delta t$ the exposure time, and $D_{\min}$ the minimum infectious dose.
The POI node's ventilation coefficient enters as an additional reduction factor
for fine suspended aerosols.

\subsection{Technology Stack}
\label{sec:tech-stack}

\subsubsection{Programming Language and Framework}
\label{sec:mesa}

The simulation is implemented in \textbf{Python} using the \textbf{Mesa}
framework~\cite{mesa2020} --- the de-facto open-source standard for agent-based
modeling in Python. Mesa provides:

\begin{itemize}
  \item a scheduling module (\texttt{RandomActivation}) --- random agent
        ordering eliminates unrealistic deterministic patterns,
  \item a spatial module (\texttt{MultiGrid}) --- grid allowing multiple agents
        per cell,
  \item \texttt{DataCollector} --- automated collection of macroscopic
        variables (S, E, I, R, D) into pandas data structures,
  \item the \texttt{mesa-geo} extension --- GIS / OpenStreetMap integration
        for realistic urban topology.
\end{itemize}

\subsubsection{Supporting Libraries}
\label{sec:libraries}

\begin{itemize}
  \item \textbf{NumPy} --- vectorised numerical operations (batch update of
        viral titres across thousands of fomites per tick),
  \item \textbf{pandas} --- results analysis and export,
  \item \textbf{Matplotlib / Seaborn} --- epidemic curve visualisation,
  \item \textbf{Streamlit} --- interactive dashboard for live parameter control
        and results inspection,
  \item \textbf{pytest} --- unit and integration tests verifying compartmental
        mechanics.
\end{itemize}

\subsubsection{Project Structure}
\label{sec:project-structure}

The project follows the \emph{src layout} convention:

\begin{lstlisting}[language=bash, caption={Planned project directory structure},
  label=lst:structure]
pandemic-simulation/
|-- pyproject.toml          # dependencies, package metadata
|-- README.md
|-- configs/                # YAML files with pathogen parameters
|   |-- pathogen_base.yaml
|   `-- scenario_lockdown.yaml
|-- data/
|   |-- input/              # demographic pyramids, GIS graphs
|   `-- output/             # DataCollector results (CSV, JSON)
|-- src/
|   `-- simulation/
|       |-- __init__.py
|       |-- agents.py       # HumanAgent, FomiteAgent
|       |-- model.py        # EpidemicModel (Mesa core)
|       |-- space.py        # POI graph, pathfinding
|       `-- transmission.py # aerosol and fomite mechanics
|-- tests/                  # pytest, unit tests
`-- scripts/                # CLI: run simulations, batch runs
\end{lstlisting}

Separating pathogen parameters into external YAML configuration files from the
core simulation logic allows switching the modelled virus (e.g.\ influenza
vs.\ a SARS variant) without modifying the simulation code itself.

%%%%%%%%%%%%%%%%%%%%%%%%%%%%%%%%%%%%%%%%%%%%%%%%%%%%%%%%%%%%%%%%%%%%%%%%%%%%%%%
\section{Stage 2: Implementation and Preliminary Verification}
\label{sec:stage2}

\emph{This section will be completed during Stage~2.}

Stage~2 covers the full implementation of all mechanisms designed in Stage~1
and preliminary verification of the model's correctness. Planned work:

\begin{itemize}
  \item full implementation of aerosol and fomite transmission mechanics as
        per the equations in Sections~\ref{sec:aerosol} and~\ref{sec:fomites},
  \item calibration of epidemiological parameters against historical data
        (COVID-19, pandemic influenza),
  \item unit and integration tests (pytest) verifying SEIRD transition
        correctness,
  \item sensitivity analysis of the model with respect to key parameters
        ($\beta$, $\sigma$, $\mu_F$, $\Theta_H$),
  \item architecture revisions informed by observations from preliminary
        simulation runs.
\end{itemize}

%%%%%%%%%%%%%%%%%%%%%%%%%%%%%%%%%%%%%%%%%%%%%%%%%%%%%%%%%%%%%%%%%%%%%%%%%%%%%%%
\section{Stage 3: Experimental Research and Final Project}
\label{sec:stage3}

\emph{This section will be completed during Stage~3.}

Stage~3 is the main research phase of the project. Planned experimental
scenarios:

\begin{enumerate}
  \item \textbf{Baseline scenario} --- pandemic without any intervention;
        observation of natural SEIRD dynamics.
  \item \textbf{Lockdown} --- closure of high-risk POI nodes (offices, shops)
        while keeping schools and healthcare facilities active.
  \item \textbf{Vaccination campaign} --- gradual rollout of a vaccine to
        selected age groups; observation of the effect on the epidemic peak.
  \item \textbf{Superspreaders} --- analysis of a scenario in which a small
        number of agents with high \texttt{contact\_rate} and low
        \texttt{hygiene\_score} account for a disproportionate number of
        infections.
  \item \textbf{Hand hygiene impact} --- comparison of populations with high
        and low $\Theta_H$ values on the dynamics of fomite-mediated
        transmission.
\end{enumerate}

Results for each scenario will be presented as epidemic curves (S, E, I, R, D
over time), infection heat maps across POI nodes, and statistical analysis of
secondary infection distributions.

%%%%%%%%%%%%%%%%%%%%%%%%%%%%%%%%%%%%%%%%%%%%%%%%%%%%%%%%%%%%%%%%%%%%%%%%%%%%%%%
\section{\SectionTitleSummary}
\label{sec:conclusions}

This document presents a comprehensive plan and incremental documentation of a
multi-scale agent-based simulation of viral pandemic spread.

The following key design decisions were made in Stage~1:

\begin{itemize}
  \item \textbf{SEIRD model} --- extension of the classical SIR model with an
        incubation state and mortality; continuous viral load as a mechanism
        linking within-host and between-host dynamics.
  \item \textbf{Agent-based architecture in Mesa} --- framework guaranteeing
        scalability, native data analysis support, and GIS integration.
  \item \textbf{Three transmission mechanisms} --- aerosol/droplets
        (distance-dependent direct transmission), fomites (indirect
        transmission via contaminated surfaces), and touch.
  \item \textbf{Heterogeneous POI space} --- multi-layered graph reflecting
        the functional structure of an urban agglomeration.
  \item \textbf{External parameter configuration} --- YAML files allow rapid
        scenario switching without modifying the core simulation code.
\end{itemize}

Sections covering Stages~2 and~3 will be filled in upon completion of each
respective project stage.

%%%%%%%%%%%%%%%%%%%%%%%%%%%%%%%%%%%%%%%%%%%%%%%%%%%%%%%%%%%%%%%%%%%%%%%%%%%%%%%
\printbibliography

%%%%%%%%%%%%%%%%%%%%%%%%%%%%%%%%%%%%%%%%%%%%%%%%%%%%%%%%%%%%%%%%%%%%%%%%%%%%%%%
\end{document}
